% 题证摘录；来源与整理边界见相同编号分题文件。

Let $G,A$ be groups and $\rho:G\to\Aut(A)$ be a homomorphism. For $a\in A$, denote $\rho(g)a$ by $g(a)$. The semidirect product $G\ltimes A$ is defined to be the product $A\times G$ with multiplication law
\[
(a_1,g_1)(a_2,g_2)=(a_1g_1(a_2),g_1g_2).
\]
Clearly, $G$ and $A$ are subgroups of $G\ltimes A$ in a natural way.

We would like to study irreducible complex representations of $G\ltimes A$. For simplicity, let us do it when $A$ is abelian. In this case, irreducible representations of $A$ are $1$-dimensional and form the character group $A^*$, which carries an action of $G$. Let $O$ be an orbit of this action, $x\in O$ a chosen element, and $G_x$ the stabilizer of $x$ in $G$. Let $U$ be an irreducible representation of $G_x$. Then we define a representation $V_{(O,U)}$ of $G\ltimes A$ as follows.

As a representation of $G$, we set
\[
V_{(O,x,U)}=\operatorname{Ind}_{G_x}^G U
=\{f:G\to U\mid f(hg)=hf(g),\ h\in G_x\}.
\]
Next, we introduce an additional action of $A$ on this space by $(af)(g)=x(g(a))f(g)$. Then it's easy to check that these two actions combine into an action of $G\ltimes A$. Also, it is clear that this representation does not really depend on the choice of $x$, in the following sense. Let $x,y\in O$, and $g\in G$ be such that $gx=y$, and let $g(U)$ be the representation of $G_y$ obtained from the representation $U$ of $G_x$ by the action of $g$. Then $V_{(O,x,U)}$ is (naturally) isomorphic to $V_{(O,y,g(U))}$. Thus we will denote $V_{(O,x,U)}$ by $V_{(O,U)}$ (remembering, however, that $x$ has been fixed).
\begin{theorem}[4.75, parts (i)--(iii)]
\begin{enumerate}[label=(\roman*)]
\item The representations $V_{(O,U)}$ are irreducible.
\item They are pairwise nonisomorphic.
\item They form a complete set of irreducible representations of $G\ltimes A$.
\end{enumerate}
\end{theorem}
\begin{proof}
(i) Let us decompose $V=V_{(O,U)}$ as an $A$-module. Then we get
\[
V=\bigoplus_{y\in O}V_y,
\]
where $V_y=\{v\in V_{(O,U)}\mid av=(y,a)v,\ a\in A\}$. (Equivalently, $V_y=\{v\in V_{(O,U)}\mid v(g)=0\text{ unless }gy=x\}$.) So if $W\subset V$ is a subrepresentation, then $W=\bigoplus_{y\in O}W_y$, where $W_y\subset V_y$. Now, $V_y$ is a representation of $G_y$, which goes to $U$ under any isomorphism $G_y\to G_x$ determined by $g\in G$ mapping $x$ to $y$. Hence, $V_y$ is irreducible over $G_y$, so $W_y=0$ or $W_y=V_y$ for each $y$. Also, if $hy=z$ then $hW_y=W_z$, so either $W_y=0$ for all $y$ or $W_y=V_y$ for all $y$, as desired.

(ii) The orbit $O$ is determined by the $A$-module structure of $V$, and the representation $U$ by the structure of $V_x$ as a $G_x$-module.

(iii) We have
\[
\begin{aligned}
\sum_{U,O}\dim V_{(U,O)}^2
&=\sum_{U,O}|O|^2(\dim U)^2
 =\sum_O|O|^2|G_x|\\
&=\sum_O|O||G/G_x||G_x|
 =|G|\sum_O|O|=|G||A^*|=|G\ltimes A|.
\end{aligned}
\]
\end{proof}
